\documentclass[11pt,a4paper]{article}

% Paquetes esenciales
\usepackage[utf8]{inputenc}
\usepackage[spanish,es-tabla]{babel}
\usepackage[margin=2.5cm]{geometry}
\usepackage{graphicx}
\usepackage{booktabs}
\usepackage{tabularx}
\usepackage{multirow}
\usepackage{amsmath}
\usepackage{hyperref}
\usepackage{xcolor}
\usepackage{titlesec}
\usepackage{float}

% Configuración de colores y enlaces
\definecolor{udacolor}{RGB}{0, 102, 51}
\hypersetup{
    colorlinks=true,
    linkcolor=udacolor,
    citecolor=blue,
    filecolor=magenta,      
    urlcolor=blue,
    pdftitle={Informe Arquitectura CubeSat UD34},
}

\titleformat{\section}{\Large\bfseries\color{udacolor}}{\thesection.}{1em}{}
\titleformat{\subsection}{\large\bfseries}{\thesubsection.}{1em}{}

\begin{document}

% ==========================================
% 1. PORTADA E IDENTIFICACIÓN
% ==========================================
\begin{titlepage}
    \centering
    \vspace*{2cm}
    {\Huge \textbf{\color{udacolor} Universidad de Antioquia}}\\[0.5cm]
    {\LARGE Facultad de Ingeniería}\\[0.5cm]
    {\Large Departamento de Ingeniería Electrónica y Telecomunicaciones}\\[2cm]
    
    {\huge \textbf{Práctica 4: Diseño Electrónico y Captura Esquemática del CubeSat UD34}}\\[2cm]
    
    {\Large \textbf{Integrantes:}}\\[0.5cm]
    {\large Juan Pablo González Blandón - juan.gonzalez57@udea.edu.co}\\[0.2cm]
    {\large Santiago Castaño Loaiza - santiago.castano2@udea.edu.co}\\[2cm]
    
    {\Large \textbf{Curso:}}\\[0.5cm]
    {\large Diseño de Productos Electrónicos}\\[2cm]
    
    \vfill
    {\large Medellín, Colombia}\\[0.2cm]
    {\large \today}
\end{titlepage}

\tableofcontents
\newpage

% ==========================================
% 1. ALCANCE Y ESPECIFICACIONES DE LA TARJETA
% ==========================================
\section{Revisión del alcance y especificaciones de la tarjeta (PCB 2: Aviónica y ADCS)}

\subsection{Identificación de Funciones Asignadas}
La \textbf{PCB 2} constituye el núcleo de procesamiento y determinación de actitud del nanosatélite UD34. De acuerdo con la arquitectura general definida en la Práctica 3, sus funciones ineludibles son:
\begin{itemize}
    \item \textbf{Procesamiento de Datos (OBC):} Ejecución del Sistema Operativo en Tiempo Real (RTOS), orquestación de rutinas de acceso directo a memoria (DMA) para el volcado de imágenes y empaquetamiento bajo el estándar CCSDS.
    \item \textbf{Almacenamiento Masivo:} Provisión de interfaces síncronas de alta velocidad para memoria volátil (búfer de ráfaga) y memoria no volátil (registro histórico de misión).
    \item \textbf{Determinación y Control de Actitud (ADCS):} Adquisición de variables inerciales y magnéticas (IMU) y síntesis de señales PWM de control hacia los puentes H que excitan las bobinas de torque (Magnetorquers).
\end{itemize}

\subsection{Trazabilidad de Requisitos y Fronteras Eléctricas}
Para garantizar el determinismo y evitar incompatibilidades con el resto del ensamble 1U, se establecen las siguientes restricciones físicas y eléctricas.

\textbf{Tensiones de Alimentación y Corrientes:}
\begin{itemize}
    \item \textbf{Lógica y Control:} $3.3\text{V}$ (dominio digital) y $5.0\text{V}$ provenientes de los convertidores de la PCB 1 (EPS). Consumo estimado nominal de $100\text{ mA}$ y picos de $250\text{ mA}$ durante el uso del DMA.
    \item \textbf{Potencia Actuadores:} Enrutamiento de $V_{BATT}$ (rango $7.4\text{V} - 8.4\text{V}$) directo desde el EPS, el cual alimentará los puentes H para inyectar picos transitorios de corriente requeridos por los Magnetorquers (estimados en $\ge 1.0\text{ A}$ a $1.5\text{ A}$ por eje).
\end{itemize}

\textbf{Asignación de Interfaces Físicas (Conectores):}
\begin{itemize}
    \item \textbf{Stack Central (Backplane PC/104):} Conector estandarizado \textit{through-hole} de apilamiento vertical que entrelaza la PCB 1 (EPS) y la PCB 3 (COMM). Agrupa los rieles de potencia, telemetría I2C (IF-05) y el enlace SPI de bajada de imágenes hacia el radio de Banda S (IF-04).
    \item \textbf{Carga Útil (Conector Mezzanine):} Conector superficial de alta densidad que apilará la cámara multiespectral (PCB 4) directamente sobre la PCB 2. Aisla el bus DCMI (IF-01) del ruido del PC/104 para transferir la ráfaga de imágenes crudas ($\ge 40\text{ MB/s}$ - Requisito ET-01).
\end{itemize}

\subsection{Resolución de Asuntos Pendientes y Modificaciones}
Durante la evaluación detallada de las interfaces, se detectó una vulnerabilidad crítica: la conmutación inductiva de alta corriente (1.5 A) de los puentes H del ADCS inducirá picos de tensión en el plano de tierra (\textit{Ground Bounce}), amenazando la integridad del bus DCMI de alta velocidad del microcontrolador y el registro de la SDRAM.

\textbf{Cambio Documentado:} Se descartan las topologías tradicionales de tierra compartida o \textit{Star Grounding}. Se integrará una arquitectura de \textbf{Aislamiento Galvánico Digital} en la PCB 2.

\textbf{Justificación Técnica y Efectos en el Sistema:}
Se implementará un aislador digital capacitivo (Texas Instruments ISO7741) y un módulo DC-DC de aislamiento galvánico de $1\text{ W}$ (Texas Instruments DCP010505BP) entre el procesador principal y la etapa del puente H (DRV8833). 
\begin{itemize}
    \item \textbf{Pros:} Anula físicamente el ruido de los Magnetorquers al desconectar la tierra digital (\texttt{DGND}) de la tierra de potencia (\texttt{PGND}), protegiendo el bus DCMI de fallos lógicos. A diferencia de un optoacoplador, el aislador capacitivo es inmune a la degradación por "Daño de Desplazamiento" radiónico en órbita baja (LEO) \cite{nasa_sst}.
    \item \textbf{Contras y Mitigación:} Exigirá sacrificar un área estimada de $2\text{ a }3\text{ cm}^2$ en la PCB 2 para alojar el conversor aislado e incrementar levemente el consumo general. Se ajustarán las reglas de ruteo para respetar el margen físico de aislamiento en la tarjeta.
\end{itemize}

% ==========================================
% 2. SELECCIÓN DE COMPONENTES
% ==========================================
```latex
% ==========================================
% 3. SELECCIÓN DE COMPONENTES
% ==========================================
\section{Selección de componentes}

La selección de componentes de la PCB 2, destinada al procesamiento de
datos y al control de actitud, considera las interfaces requeridas, los
rangos de alimentación, el consumo, el encapsulado y la disponibilidad
comercial. Se priorizan componentes COTS (\textit{Commercial Off-The-Shelf})
por su integración y disponibilidad; sin embargo, su utilización no
implica por sí misma una calificación para operación espacial. Su
compatibilidad con el entorno de la misión deberá evaluarse según los
requisitos térmicos, mecánicos y de radiación aplicables.

\subsection{Componentes principales}

\begin{table}[htbp]
\centering
\renewcommand{\arraystretch}{1.35}
\small
\begin{tabularx}{\textwidth}{
|>{\raggedright\arraybackslash}p{2.6cm}
|>{\raggedright\arraybackslash}p{3.0cm}
|X|}
\hline
\textbf{Función} & \textbf{Referencia seleccionada} &
\textbf{Justificación, condiciones y aspectos por verificar} \\
\hline

Procesamiento central &
STM32H753IIT6 &
Microcontrolador basado en ARM Cortex-M7, con frecuencia máxima de
480 MHz, 2 MB de memoria Flash y 1 MB de RAM interna. Integra
interfaces como SPI, I\textsuperscript{2}C, UART y FMC, útiles para
la conexión de periféricos y memorias externas. La frecuencia de
operación definitiva debe justificarse según la carga de procesamiento,
el consumo, el reloj seleccionado y los límites eléctricos del
dispositivo. Se deben dimensionar los circuitos de alimentación,
reinicio, arranque y depuración a partir de la hoja de datos. \\
\hline

Memoria volátil externa &
ISSI IS42S16400J-7TL &
SDRAM síncrona de 64 Mbit, equivalente a 8 MiB, con bus de datos de
16 bits. Su interfaz paralela permite utilizar el controlador FMC del
microcontrolador. La temporización debe configurarse a partir de los
parámetros AC de la referencia exacta y de la frecuencia del reloj
de memoria. El sufijo de velocidad no debe interpretarse por sí solo
como el tiempo de acceso. Se deben verificar la temperatura de
operación, los requisitos de refresco y los niveles de alimentación. \\
\hline

Memoria no volátil &
Kingston EMMC04G-M627-A01 &
Memoria eMMC de 4 GB con controlador interno, apropiada para
almacenamiento masivo. La referencia admite una alimentación VCC
de 2,7--3,6 V y distintas opciones de VCCQ, según el modo de
operación. Para el diseño se considera inicialmente VCC de 3,3 V
y VCCQ de 1,8 V, sujeto a verificar la compatibilidad con la interfaz
y la configuración del host. No se asume tolerancia a radiación por
el hecho de incorporar mecanismos internos de gestión de errores. \\
\hline

Unidad de medición inercial &
TDK InvenSense ICM-20948 &
Integra acelerómetro, giroscopio y magnetómetro en un dispositivo
MEMS de nueve ejes, con comunicación digital SPI o I\textsuperscript{2}C.
Su integración reduce el número de componentes externos y facilita
la estimación de orientación. El fabricante identifica esta referencia
como EOL (\textit{End of Life}), por lo que su disponibilidad representa
un riesgo para la fabricación y el mantenimiento. Se conserva como
referencia inicial, teniendo en cuenta dicho riesgo. \\
\hline

Control de actuadores &
Texas Instruments DRV8833 &
Integra dos puentes H para controlar cargas inductivas. Su rango de
alimentación de motor es de 2,7--10,8 V. La corriente utilizable depende
del encapsulado, la temperatura y el diseño térmico. La corriente
objetivo de cada bobina debe determinarse antes de confirmar que
este driver sea adecuado. Para tres ejes de actuación independientes
se debe verificar la cantidad de puentes H necesarios y la asignación
de señales de control. \\
\hline

Aislamiento digital &
Texas Instruments ISO7741 &
Aislador digital de cuatro canales, con tres canales en una dirección
y uno en la dirección opuesta. Admite niveles lógicos dentro del
rango especificado por el fabricante, pero no integra una fuente de
alimentación aislada. La selección final depende del mapa de señales,
del número de canales requeridos y de los requisitos de aislamiento
del diseño. \\
\hline

Conversión DC-DC aislada &
Texas Instruments DCP010505BP &
Convertidor DC-DC aislado de salida nominal de 5 V y potencia nominal
de hasta 1 W para la variante correspondiente. Su capacidad debe
verificarse frente a la carga real del dominio aislado. No se considera
como fuente de potencia para las bobinas, debido a que su capacidad
de salida puede resultar insuficiente para la corriente requerida por
los actuadores. \\
\hline
\end{tabularx}
\caption{Componentes principales considerados para la PCB 2.}
\label{tab:componentes_pcb2}
\end{table}

\subsection{Criterios complementarios de selección}

La selección de los componentes pasivos y auxiliares se realizará
conforme a las hojas de datos y a las condiciones de operación
previstas para cada circuito:

\begin{itemize}
    \item \textbf{Desacoplamiento:} se consideran condensadores
    cerámicos próximos a los pines de alimentación de los circuitos
    integrados. Sus valores y cantidades se definirán según las
    recomendaciones del fabricante y los transitorios previstos.
    La capacidad de reserva de la etapa de potencia se dimensionará
    por separado.

    \item \textbf{Relojes:} se considera inicialmente un cristal HSE
    de 25 MHz y, si los requisitos del sistema lo justifican, un
    cristal LSE de 32,768 kHz para el RTC. La selección definitiva
    depende de la capacitancia de carga, la tolerancia, el rango
    térmico y las recomendaciones del oscilador del microcontrolador.

    \item \textbf{Buses digitales:} las resistencias de polarización
    de I\textsuperscript{2}C se calculan a partir de la tensión del
    bus, la corriente de hundimiento admisible y la capacitancia total
    de las líneas. Las resistencias serie para señales rápidas se
    definirán según las características eléctricas y el trazado de
    la PCB, evitando adoptar valores genéricos sin verificación.

    \item \textbf{Configuración y depuración:} se prevén conexiones
    para programación, reinicio y selección del modo de arranque del
    STM32. Los pines no utilizados se tratarán según las indicaciones
    del fabricante, evitando entradas flotantes cuando no estén
    permitidas.

    \item \textbf{Compatibilidad eléctrica:} se comprobarán las
    tensiones de alimentación y los niveles lógicos de cada interfaz,
    especialmente la compatibilidad entre VCCQ de la eMMC y el
    microcontrolador, así como las conexiones entre los dominios
    aislados.
\end{itemize}


% ==========================================
% 4. DISEÑO Y DIMENSIONAMIENTO
% ==========================================
\section{Diseño y dimensionamiento de circuitos}

El dimensionamiento de los circuitos de soporte busca establecer
valores compatibles con las condiciones de operación recomendadas
por los fabricantes y con los requisitos funcionales de la PCB 2.
Los cálculos se plantean a partir de los parámetros eléctricos
disponibles y de las hipótesis iniciales de diseño.

\subsection{Alimentación y dominios eléctricos}

Se considera como punto de partida una batería de dos celdas de
ion-litio conectadas en serie, con una tensión nominal de 7,4 V y
una tensión máxima de 8,4 V. Para el diseño preliminar se adopta
un intervalo de entrada de 6,0--8,4 V, cuya compatibilidad con la
batería y su límite de descarga deberá mantenerse como condición
de diseño.

A partir de este bus se contempla un dominio lógico de 3,3 V para
el microcontrolador y los periféricos compatibles. Las memorias y
los sensores se alimentarán de acuerdo con sus rangos especificados,
utilizando regulación adicional cuando sea necesaria.

El dominio de los actuadores deberá mantenerse galvánicamente
aislado del dominio lógico. Para ello, se requiere aislamiento de
las señales de control y una fuente de alimentación aislada
dimensionada según la potencia de la carga. El ISO7741 se utiliza
para el aislamiento digital, mientras que el convertidor DC-DC se
seleccionará de acuerdo con los requisitos de corriente y tensión
del dominio aislado.

El DCP010505BP no se asigna como fuente de alimentación directa
de los magnetorquers, debido a que su potencia nominal de salida
es de hasta 1 W. La fuente definitiva para los actuadores deberá
cubrir la demanda eléctrica calculada de las bobinas y las pérdidas
del convertidor.

\subsection{Dimensionamiento de la red de reloj HSE}

Para dimensionar los condensadores de carga de un cristal externo,
se utiliza la relación:

\[
C_L =
\frac{C_1 C_2}{C_1+C_2}+C_{\mathrm{stray}},
\]

donde \(C_L\) es la capacitancia de carga especificada por el
fabricante del cristal, \(C_1\) y \(C_2\) son los condensadores
externos y \(C_{\mathrm{stray}}\) representa la capacitancia parásita
del circuito.

Si se utiliza una capacitancia externa simétrica,
\(C_1=C_2=C_{\mathrm{ext}}\), entonces:

\[
C_{\mathrm{ext}}=2(C_L-C_{\mathrm{stray}}).
\]

Como cálculo preliminar, considerando una capacitancia de carga
del cristal de \(C_L=8\,\mathrm{pF}\) y una capacitancia parásita
estimada de \(C_{\mathrm{stray}}=2\,\mathrm{pF}\), se obtiene:

\[
C_{\mathrm{ext}}=2(8-2)\,\mathrm{pF}
=12\,\mathrm{pF}.
\]

Por tanto, se considera inicialmente un condensador de \(12\,\mathrm{pF}\)
en cada terminal del cristal, con dieléctrico C0G/NP0. Esta selección
queda condicionada a confirmar la referencia exacta del cristal,
su capacitancia de carga y las recomendaciones del oscilador HSE
del STM32H753.

\subsection{Resistencias de polarización del bus I\textsuperscript{2}C}

El valor de las resistencias de polarización depende de la tensión
del bus, la corriente de hundimiento permitida por los dispositivos
y la capacitancia total de las líneas. Para un cálculo preliminar,
se emplean:

\[
R_{p(\min)}=\frac{V_{CC}-V_{OL}}{I_{OL}}
\]

y

\[
R_{p(\max)}=\frac{t_r}{0.8473\,C_b},
\]

donde \(V_{OL}\) es la tensión de nivel bajo, \(I_{OL}\) la corriente
de hundimiento admisible, \(t_r\) el tiempo de subida máximo y
\(C_b\) la capacitancia total del bus.

Considerando \(V_{CC}=3.3\,\mathrm{V}\), \(V_{OL}=0.4\,\mathrm{V}\),
\(I_{OL}=3\,\mathrm{mA}\), \(t_r=300\,\mathrm{ns}\) y
\(C_b=50\,\mathrm{pF}\), se obtiene:

\[
R_{p(\min)}=
\frac{3.3-0.4}{3\,\mathrm{mA}}
\approx 967\,\Omega
\]

\[
R_{p(\max)}=
\frac{300\,\mathrm{ns}}
{0.8473(50\,\mathrm{pF})}
\approx 7.08\,\mathrm{k}\Omega.
\]

Bajo estas condiciones, se propone inicialmente una resistencia
de \(4.7\,\mathrm{k}\Omega\), que se encuentra dentro del intervalo
calculado. Su aceptación final depende de confirmar la capacitancia
real del bus, el modo de operación y los límites eléctricos de
todos los dispositivos conectados.

\subsection{Dimensionamiento de los magnetorquers y la etapa de potencia}

Para iniciar el dimensionamiento de los magnetorquers se adopta
como objetivo preliminar un momento dipolar de
\(0.07\,\mathrm{A\,m^2}\) por eje, tomado como referencia de diseño
para un CubeSat de clase 1U y no como un requisito oficial confirmado
de la misión UD34.

El momento dipolar de una bobina puede aproximarse mediante:

\[
m=NIA_{\mathrm{ef}},
\]

donde \(N\) es el número de espiras, \(I\) la corriente de la bobina
y \(A_{\mathrm{ef}}\) el área efectiva por espira. El par magnético
generado en el campo magnético terrestre se expresa como:

\[
\boldsymbol{\tau}=\boldsymbol{m}\times\boldsymbol{B},
\]

donde \(\boldsymbol{B}\) es el campo magnético local.

Para una bobina implementada mediante pistas de cobre en la PCB,
la resistencia eléctrica inicial se estima con:

\[
R=\rho\frac{\ell}{wt},
\]

donde \(\rho\) es la resistividad del cobre, \(\ell\) la longitud
total de la pista, \(w\) su ancho y \(t\) el espesor del cobre.
La potencia disipada por efecto Joule y la tensión resistiva en
régimen estacionario se calculan mediante:

\[
P_{\mathrm{Cu}}=I^2R,
\qquad
V_R=IR.
\]

Se propone una configuración de tres ejes magnéticos independientes.
Para lograrlo, la geometría debe permitir componentes del momento
dipolar en tres direcciones; tres bobinas coplanares no generan
por sí solas tres ejes independientes.

El objetivo dipolar adoptado permite establecer una referencia
inicial, pero la corriente, la resistencia, la inductancia y la
potencia de cada bobina dependen de su geometría física, número
de espiras, ancho y espesor de las pistas. Por esta razón, dichos
parámetros deberán obtenerse a partir de la geometría definitiva
de la PCB antes de confirmar la capacidad del DRV8833.

La corriente admisible del driver deberá verificarse para el
encapsulado seleccionado, considerando la resistencia de conducción,
la disipación térmica y las condiciones de operación. Si se emplea
una resistencia de sensado, su valor se calculará a partir del
umbral y la función de sensado o regulación de corriente indicada
por el fabricante. No se fija un valor definitivo de \textit{shunt}
sin haber establecido previamente la corriente objetivo de la bobina.

\subsection{Presupuesto de consumo de la PCB 2}

El consumo de cada dominio se calcula a partir de la tensión de
alimentación y la corriente correspondiente al estado de operación:

\[
P_i=V_iI_i.
\]

Cuando se utiliza un convertidor, la potencia requerida en su
entrada depende de la potencia suministrada a la carga y de su
eficiencia:

\[
P_{\mathrm{entrada}}=\frac{P_{\mathrm{carga}}}{\eta},
\]

donde \(\eta\) representa la eficiencia del convertidor.

Para el presupuesto de la PCB 2 se distinguen los siguientes estados:

\begin{itemize}
    \item \textbf{Estado nominal:} microcontrolador ejecutando el
    software de control, sensores activos, memorias en operación
    habitual e interfaces digitales habilitadas según el modo de
    misión.

    \item \textbf{Estado de máxima demanda previsto:} procesamiento
    y transferencias de memoria activas, operaciones de escritura
    en la eMMC, comunicaciones habilitadas y actuación de los
    magnetorquers según la estrategia de control.

    \item \textbf{Estado de reposo:} periféricos deshabilitados o
    en bajo consumo, según la configuración de cada componente.
    La corriente de reposo de los convertidores y del dominio
    aislado también debe incluirse cuando corresponda.
\end{itemize}

Las corrientes nominales y máximas deberán obtenerse de las hojas
de datos para las condiciones correspondientes de tensión,
frecuencia, temperatura y actividad. La corriente de los actuadores
se calculará a partir de las bobinas dimensionadas y no únicamente
de la corriente máxima admisible del driver. Asimismo, las
corrientes máximas de los componentes no se sumarán como si
ocurrieran simultáneamente sin verificar que dicho estado sea
posible en la operación prevista.

Por lo anterior, no se asigna todavía un valor numérico definitivo
a la potencia total de la PCB 2. El presupuesto deberá integrar
las corrientes verificadas de los circuitos digitales, las pérdidas
de conversión y la potencia calculada de los magnetorquers.

\subsection{Configuración de arranque, depuración y tratamiento de pines}

Para garantizar la fiabilidad del sistema embebido y facilitar las pruebas terrestres, se definen las siguientes políticas de diseño a nivel esquemático:

\begin{itemize}
    \item \textbf{Condiciones de arranque (\textit{Boot}):} Para evitar que un evento singular (SEU) altere la secuencia de inicio, el pin \texttt{BOOT0} del STM32H753 se fijará a tierra (\texttt{DGND}) mediante una resistencia \textit{pull-down} de $10\,\mathrm{k\Omega}$. Esto asegura un arranque determinista exclusivamente desde la memoria Flash principal.
    \item \textbf{Programación y Depuración:} Se implementará una interfaz \textbf{SWD (Serial Wire Debug)} de 4 pines (\texttt{SWDIO}, \texttt{SWCLK}, \texttt{VDD}, \texttt{GND}). Esta topología minimiza el uso de pines respecto al estándar JTAG, reservando puertos GPIO para el enrutamiento de la carga útil.
    \item \textbf{Tratamiento de pines no utilizados:} Los pines sobrantes del microcontrolador permanecerán físicamente desconectados en la PCB para optimizar el área de enrutamiento. Su mitigación contra acoplamientos electromagnéticos se delegará al \textit{firmware} de inicialización (configurándolos como entradas con \textit{pull-down} o salidas analógicas).
    \item \textbf{Puntos de prueba (\textit{Test Points}):} Se preverán superficies de cobre expuestas (Pads tipo \textit{Test Point}) en señales críticas para la validación funcional: niveles de voltaje aislados (\texttt{3.3V\_Iso}), línea de reinicio de hardware (\texttt{NRST}), voltajes de capacitores internos del núcleo (\texttt{VCAP}) y las señales PWM de excitación del puente H.
\end{itemize}

\subsection*{Nota sobre el carácter preliminar del diseño}

Los componentes seleccionados, los valores calculados y las
estimaciones presentadas corresponden a una etapa previa al diseño
físico de la PCB, la definición definitiva de sus dimensiones y
capas, y el enrutamiento de las pistas. Por tanto, constituyen una
base inicial de diseño y pueden estar sujetos a modificaciones
durante la captura esquemática, la simulación, el análisis de
integridad de señal y el diseño del circuito impreso.

En particular, el dimensionamiento definitivo de los magnetorquers
y el cierre del presupuesto de consumo dependen de la geometría
final de las bobinas, las características físicas de las pistas de
cobre y la confirmación del momento dipolar requerido por la misión.
Estos parámetros deberán verificarse antes de considerar cerrados
el diseño eléctrico y la selección de la etapa de potencia.
```


% ==========================================
% REFERENCIAS
% ==========================================
\begin{thebibliography}{99}
\bibitem{ecss} European Cooperation for Space Standardization (ECSS), \textit{ECSS-E-ST-20-20C: Electrical design and interface requirements for power supply}, 2022.
\bibitem{nasa_sst} NASA, \textit{State-of-the-Art Small Spacecraft Technology: Communications and Hardware Systems}, 2026.
\end{thebibliography}

\end{document}